\section{Domain-agnostic reconstruction framework}
\label{sec:framework}

\subsection{Problem formulation}
\label{sec:problem}

Let $X = \{x_0, \ldots, x_n\}$ be a time-ordered trajectory of an aircraft or
a vessel. Each state $x_i = (y_i, k_i, a_i, \mathrm{time}_i)$ contains a
horizontal position $y_i = (\phi_i, \lambda_i)$, kinematic fields $k_i$ (for
example ground speed and track, or SOG and COG) and the reading $a_i$ of a
trusted auxiliary signal. The trusted signal must be \emph{physically}
independent of the satellite receiver, not merely reported in a separate
field. An external detector raises an alarm at index $t$. From then on,
every satellite-derived field ($y_i$ and the satellite-derived part of
$k_i$, for $i > t$) is withheld from the reconstruction, while $a_i$ remains
available. The task is to produce
\begin{equation}
\hat{Y}_{t+1:r} = \{\hat{y}_i\}_{i=t+1}^{r}
\label{eq:task}
\end{equation}
until trusted reporting resumes, or the trip ends, at index $r$. The
available information is the verified prefix $X_{0:t}$, the trusted signal
$a_{0:r}$, environmental context obtained independently of the asset, and
historical trajectories of the same domain. This is not denoising, because
corrupted observations are dropped rather than filtered. It is not
interpolation either, because the end state is unknown. We evaluate it as
forward-only reconstruction against the true, uncorrupted positions.

\subsection{Shared engine and domain adapter}

We factor the method into a shared engine and a domain adapter. The
\emph{engine} comprises a frozen text-embedding model $g$, a frozen LLM $f$,
an inner-product vector index, the iterative propose--validate--update loop
of Algorithm~\ref{alg:reconstruction}, the feasibility guard and a fixed
output schema (one latitude/longitude pair per step). No component of the
engine is trained, fine-tuned or adjusted for a domain. A domain adapter is
the tuple
\begin{equation}
\mathcal{D} = \langle \mathcal{S}, a, \mathcal{P}, \pi, \tau, \mathcal{E}, \mathcal{M}, \kappa \rangle,
\label{eq:adapter}
\end{equation}
whose elements are:
\begin{itemize}
\item $\mathcal{S}$, the \emph{state schema}, which lists the fields of a
surveillance report and marks each as satellite-derived (untrusted after
onset) or trusted;
\item $a$, the \emph{trusted auxiliary signal};
\item $\mathcal{P}$ and $\pi$, the \emph{phase taxonomy} and a
phase-assignment rule that uses only trusted information;
\item $\tau$, the \emph{verbaliser}, a prompt template that turns a
trajectory segment, the trusted signal and the environmental context into a
short textual state description;
\item $\mathcal{E}$, the \emph{environmental context provider}, queried at
the last accepted (trusted or reconstructed) position and time;
\item $\mathcal{M} = \{\mathcal{M}_p\}_{p \in \mathcal{P}}$, the
\emph{phase-indexed memory}, built by applying $\tau$ and $g$ to historical
trajectories of the domain;
\item $\kappa = (\Delta t, \{\alpha_p\}, \mathrm{SPEEDRULE})$, the
\emph{kinematic constants}: the sampling interval, per-phase tolerance
factors and the rule that updates the reference speed during a blackout.
\end{itemize}
Three of these elements are semantic choices that a domain expert makes
once: the trusted signal, the phase taxonomy and the vocabulary of the
verbaliser. The remaining elements are data artefacts produced from the
domain corpus by procedures that are identical in both domains
(Sections~\ref{sec:offline-memory} and \ref{sec:calibration}). Porting the
method to a new domain therefore requires a prompt template, an index and a
calibration run. It requires no training run.

\begin{algorithm}[t]
\caption{Online reconstruction with the shared engine and a domain adapter $\mathcal{D}$}
\label{alg:reconstruction}
\begin{algorithmic}[1]
\Require verified prefix $X_{0:t}$; trusted signal $a_{0:r}$; adapter
$\mathcal{D} = \langle \mathcal{S}, a, \mathcal{P}, \pi, \tau, \mathcal{E}, \mathcal{M}, \kappa \rangle$
with $\kappa = (\Delta t, \{\alpha_p\}, \mathrm{SPEEDRULE})$; frozen LLM $f$;
frozen encoder $g$; retrieval size $K$
\Ensure reconstructed positions $\hat{Y}_{t+1:r}$
\State $p \gets \pi(X_{0:t}, a_{0:t})$ \Comment{phase from trusted information only}
\State $(\hat{y}_t, v) \gets (y_t, v_t)$ \Comment{last verified position and speed}
\State $q \gets \tau(X_{0:t}, a_{0:t}, (\hat{y}_t, t))$ \Comment{textual state of the prefix}
\For{$i = t+1, \ldots, r$}
  \State $C \gets \mathrm{TopK}(g(q), \mathcal{M}_p, K)$ \Comment{cosine similarity}
  \State $\tilde{y}_i \gets f(q, C, a_i, (\hat{y}_{i-1}, i))$ \Comment{fixed schema: latitude, longitude}
  \State $d_{\max} \gets \alpha_p \, v \, \Delta t$
  \If{$d_{\mathrm{hav}}(\tilde{y}_i, \hat{y}_{i-1}) > d_{\max}$}
    \State $\tilde{y}_i \gets \mathrm{Project}(\tilde{y}_i, \hat{y}_{i-1}, d_{\max})$; raise low-confidence flag
  \EndIf
  \State $\hat{y}_i \gets \tilde{y}_i$
  \State $v \gets \mathrm{SPEEDRULE}(v, \hat{y}_{i-1}, \hat{y}_i, i)$ \Comment{domain-specific (Section~\ref{sec:instantiation})}
  \State $q \gets \tau(X_{0:t} \cup \hat{Y}_{t+1:i}, a_{0:i}, (\hat{y}_i, i))$ \Comment{feedback: re-query from reconstructed prefix}
\EndFor
\State \Return $\hat{Y}_{t+1:r} = \{\hat{y}_i\}_{i=t+1}^{r}$
\end{algorithmic}
\end{algorithm}

\subsection{Offline memory construction}
\label{sec:offline-memory}

Historical trajectories are cleaned (duplicates removed, timestamps made
monotonic), resampled to the adapter's $\Delta t$ when the native sampling
is irregular, and segmented into phases. Each segment is verbalised by
$\tau$ into a description of its spatial evolution, kinematic trend, phase,
elapsed time in phase and aligned environmental context. Timestamps are
expressed as time from the start of the phase so that segments recorded on
different days are comparable. Descriptions are embedded by $g$,
L2-normalised and stored in a separate inner-product index per phase (and,
in the maritime adapter, per corridor). Each index entry keeps a pointer to
the numeric continuation of its segment, which is shown to the LLM at
inference time. Separating indices by phase prevents incompatible motion
regimes, such as berthing and open-channel transit, from being retrieved
together. New data are added by appending entries, without re-fitting any
model.

\subsection{Online retrieval-conditioned reconstruction}

At onset, the phase is determined by $\pi$ from trusted information, and the
verified prefix is verbalised into a query. At each step the engine
retrieves the $K$ most similar segments from $\mathcal{M}_p$ by cosine
similarity. The LLM then receives the query, the retrieved segments with
their continuations, the trusted reading $a_i$ and the environmental
context, and returns the next position in the fixed schema. After the
position passes the guard, the query is rebuilt from the extended (verified
+ reconstructed) prefix, so the retrieved neighbours evolve as the
reconstruction progresses (Algorithm~\ref{alg:reconstruction}). The LLM thus
acts as a retrieval-conditioned \emph{proposal generator}: it adapts and
blends the retrieved precedents to the current context, but it never has the
final word on feasibility.

\subsection{Kinematic feasibility guard and its calibration procedure}
\label{sec:calibration}

Each proposal $\tilde{y}_i$ must lie within a disc centred at the last
accepted position, with radius
\begin{equation}
d_{\max} = \alpha_p \, v \, \Delta t,
\label{eq:guard}
\end{equation}
where $v$ is the current reference speed and $\alpha_p$ is a phase
tolerance. Proposals outside the disc are projected onto its boundary along
the bearing from the previous position, and the step is flagged as low
confidence. The guard therefore turns any LLM error, including a
hallucinated position, into a bounded displacement.

The tolerance is not tuned on reconstruction accuracy. It is set by a fixed,
accuracy-independent \emph{calibration procedure}. On the domain's training
trajectories we compute the step displacement ratio
\begin{equation}
\rho_i = \frac{d_{\mathrm{hav}}(y_{i-1}, y_i)}{v_{i-1} \, \Delta t},
\label{eq:rho}
\end{equation}
excluding steps with near-zero reference speed, and we set $\alpha$ to a
high quantile of $\rho$ (the 99th percentile in this study). For aircraft
this procedure gives $Q_{0.99}(\rho) = 1.18$ and hence $\alpha = 1.2$
\citep{felendler2026trajrag}. Section~\ref{sec:maritime-calibration} applies
the same procedure to vessels and shows why the procedure, rather than the
constant, is what transfers between domains.

\subsection{Real-time dual-domain system architecture}
\label{sec:architecture}

Fig.~\ref{fig:architecture} shows how the framework is deployed as one
service for both surveillance domains. ADS-B and AIS streams feed an
external spoofing detector, which this paper treats as a given interface.
When an alarm is raised for an asset, a trust gate stops forwarding that
asset's satellite-derived fields. The message type selects the domain
adapter. The shared engine then produces one reconstructed position per
reporting interval until the detector clears the alarm or the trip ends.
Each reconstructed position carries a confidence flag that records whether
the guard had to clamp the proposal, so a controller or VTS operator sees
the reconstruction together with an explicit indication of how far it
relies on the physical bound rather than on precedent. Because adapters
contain only prompts, indices and constants, a new domain or corridor is
added by deploying a new adapter. The engine, and any accreditation
evidence gathered for it, stays unchanged. Memories are refreshed by
appending completed trips, for example once a day.
Section~\ref{sec:latency} checks the latency of this loop against the
reporting interval of each channel.

\begin{figure}[t]
\centering
\resizebox{\linewidth}{!}{%
\begin{tikzpicture}[
  node distance=10mm and 8mm,
  box/.style={draw, rounded corners, align=center, minimum height=9mm, font=\footnotesize, inner sep=3pt},
  adapter/.style={box, fill=orange!15},
  engine/.style={box, fill=blue!10},
  io/.style={box, fill=gray!10},
  >=Stealth
]
\node[io] (adsb) {ADS-B stream\\(aircraft)};
\node[io, below=6mm of adsb] (ais) {AIS stream\\(vessels)};
\node[io, right=14mm of adsb, yshift=-6mm] (detector) {External\\spoofing detector};
\node[io, right=10mm of detector] (gate) {Trust gate: drop\\satellite fields\\after onset};
\node[adapter, right=10mm of gate, yshift=7mm] (adaptA) {Aviation adapter\\$\tau$, baro.\ alt., $\kappa$};
\node[adapter, below=4mm of adaptA] (adaptM) {Maritime adapter\\$\tau$, gyro heading, $\kappa$};
\node[adapter, right=8mm of adaptA, yshift=-7mm, minimum width=24mm] (memory) {Phase-indexed\\memory $\mathcal{M}$};
\node[engine, right=10mm of memory, minimum width=26mm] (retrieval) {Phase-aware retrieval:\\returns top-$K$ precedents};
\node[engine, right=4mm of retrieval] (llm) {Frozen LLM\\next-position proposal};
\node[engine, right=4mm of llm] (guard) {Kinematic\\feasibility guard};
\node[io, right=10mm of guard] (out) {Reconstructed track\\+ confidence flag\\$\to$ ATC / VTS};
\draw[->] (adsb) -- (detector);
\draw[->] (ais) -- (detector);
\draw[->] (detector) -- node[above,font=\tiny]{alarm} (gate);
\draw[->] (gate) -- node[above,font=\tiny]{ADS-B} (adaptA);
\draw[->] (gate) -- node[below,font=\tiny]{AIS} (adaptM);
\draw[->] (adaptA) -- (memory);
\draw[->] (adaptM) -- (memory);
\draw[->] (memory) -- (retrieval);
\draw[->] (retrieval) -- (llm);
\draw[->] (llm) -- (guard);
\draw[->] (guard) -- (out);
\draw[->] (guard.south) to[out=270,in=270,looseness=1.3] node[below,font=\tiny]{feedback: re-query, extended prefix} (retrieval.south);
\draw[->, dashed] (gate.north) to[out=90,in=90,looseness=0.5] node[above,font=\tiny]{live trusted reading $a_i$, every step} (llm.north);
\end{tikzpicture}%
}
\caption{Architecture of the real-time, domain-agnostic system. Both
surveillance streams pass through the same detector interface and trust
gate. A domain adapter (orange) supplies the domain-specific parts: the
textual-state prompt template, the trusted non-satellite signal, the
kinematic constants, and the phase-indexed memory $\mathcal{M}$ built from
that domain's history -- shown as a separate node from the shared engine's
retrieval step (blue), which searches $\mathcal{M}$ and returns only the
top-$K$ precedents for the current query, not the memory itself. The dashed
arrow shows the trusted reading $a_i$ (heading or barometric altitude)
entering the loop live at every step, independently of the one-shot
adapter routing above it. The engine asks a frozen large language model for
the next position and accepts it only inside the kinematic envelope; the
accepted position extends the verified prefix, and the query is rebuilt
from it before the next retrieval. ATC: air traffic control; VTS: vessel
traffic service.}
\label{fig:architecture}
\end{figure}
